\paragraph{生産者物価テイラールール（TR-PPI）}
\label{sec:chapter5_a_shock_policies_tr_ppi}
TR-PPIの式については
\ref{sec:chapter5_beta_shock_policies_tr_ppi}
を参照のこと。
\begin{enumerate}
    \item
    \label{itm:chapter5_a_shock_policies_tr_ppi_c_H_W}
        \(c_t^{H\to W}\)：
        総消費指数\(c_t^{H\to W}\)の図\ref{fig:chapter5_a_shock_graphs_c_H_W}を見ると、TR-PPIの軌道は生産者物価インフレ目標（IT-PPI）や消費者物価テイラールール（TR-CPI）などと同様に、ショック直後の初期において極端な大暴落を引き起こしていることがわかる。
        総消費水準目標（CLT）や名目総消費水準目標（NCLT）が落ち込みをほぼゼロに抑え込んでいるのとは対照的である。
        供給ショックに伴う生産者物価のインフレ圧力に対して、テイラールールに基づく機械的な金利の引き上げ（金融引き締め）で対応してしまうため、実体経済への強い下押し圧力が働き、消費の深刻な減少を招いている。
        消費の低下は厚生にマイナスに働くため、この前半における深刻な消費の落ち込みがTR-PPIの厚生を全体の中で下位クラスへと低迷させる大きな要因となっている。
        また深く落ち込んだ状態から極めて長い時間をかけて定常状態へ回帰するだけの軌道となっているため、消費増加による厚生の改善効果も得られていない。
    \item
    \label{itm:chapter5_a_shock_policies_tr_ppi_y_H}
        \(y_t^H\)：
        生産\(y_t^H\)の図\ref{fig:chapter5_a_shock_graphs_y_H}を見ると、TR-PPIの軌道は総消費指数\(c_t^{H\to W}\)と同様に、生産者物価インフレ目標（IT-PPI）や消費者物価テイラールール（TR-CPI）と極めて類似した大暴落を見せていることがわかる。
        前述の通り、インフレ圧力に対する機械的な金融引き締めが実体経済をさらに圧迫するため、ショック直後の前半においてTR-PPIはマイナス方向へ極めて大きく落ち込んでいる。
        生産の低下は労働の減少を意味するため、この前半部分においては労働の非効用を減らして厚生にプラスに働く効果が他の政策よりも大きくなっている。
        しかしながら、前述した深刻な消費の大暴落によるマイナス効果や、後述する価格分散の増大によるマイナス効果を補うほどの改善には至っておらず、結果として全体的な厚生の評価を低迷させる結果となっている。
        また深く落ち込んだ状態から極めて長い時間をかけて定常状態へ回帰する軌道を描いている。
    \item
    \label{itm:chapter5_a_shock_policies_tr_ppi_pi_H}
        \(\pi_t^H\)：
        グロス・インフレ率\(\pi_t^H\)の図\ref{fig:chapter5_a_shock_graphs_pi_H}を見ると、TR-PPIの軌道は生産者物価インフレ目標（IT-PPI）のように完全に平坦な推移を維持できているわけではなく、ショック直後の初期において供給ショックの直接的な影響を受け、消費者物価目標（CPLTやIT-CPI）と同程度のプラス方向（インフレ）への鋭いスパイクを許容してしまっていることがわかる。
        しかし、テイラールールに従って生産者物価の変動に機械的な金融引き締めで対応するため、その後は国内生産財のインフレ率の変動を強く抑制しゼロ近傍に引き戻している。この中盤以降のインフレ率の安定化が次項で述べる価格分散の過度な蓄積を抑える要因となっている一方で、前述した深刻な実体経済の大暴落を招く代償を払っていると言える。
    \item
    \label{itm:chapter5_a_shock_policies_tr_ppi_Delta_t_minus_1_H}
        \(\Delta_{t-1}^H\)：
        価格分散\(\Delta_{t-1}^H\)の図\ref{fig:chapter5_a_shock_graphs_Delta_H}を見ると、TR-PPIの定常状態からの乖離は、生産者物価インフレ目標（IT-PPI）や生産者物価水準目標（PPLT）と同様に全期間を通じてほぼゼロの水準を維持していることがわかる。
        総消費水準目標（CLT）や生産水準目標（OLT）などが極めて大きく増大しているのとは対照的である。
        これは前述のインフレ率の安定化により、価格改定に伴う相対価格のばらつきが抑制された結果であり、この面においては生産効率の悪化による厚生へのマイナス効果を防いでいる。
        しかしながら、供給ショック（\(a\)ショック）においては価格分散の増大を抑制することのプラス効果以上に、実体経済（消費や生産）の大暴落によるマイナス効果が極めて絶大である。そのため、価格分散の発生を抑え込んだにもかかわらず、初期の深刻な不況を防げなかったことが致命傷となり、TR-PPIの厚生は全体の中で下位クラスへと低迷する結果となっている。
\end{enumerate}